

%%%%%
% Standard packages
%%%%%


% \usepackage[compress]{cite} % numbered ordering of multiple citations, set to compress ranges
\usepackage{amsmath,amsfonts}
\usepackage{algorithmic}
\usepackage{graphicx}
\usepackage{textcomp}
\usepackage{booktabs} % For formal tables \toprule
\usepackage{xspace} % Put an \xspace within numeric macros
\usepackage[normalem]{ulem}
\usepackage{makecell}
% \usepackage{tcolorbox}
\usepackage{siunitx}
\usepackage[vskip=1em,font=itshape,leftmargin=2em,rightmargin=2em]{quoting}
% \usepackage{svg}
\usepackage{listings}
\usepackage{subfigure}
\usepackage{url}
\def\UrlBreaks{\do\/\do-\do_} % Defining where it is okay to wrap a url to a new line. For bibliography.
%\useunder{\uline}{\ul}{}
% \usepackage{filecontents,lipsum}
% \usepackage[noadjust]{cite}


% \usepackage{pifont}% http://ctan.org/pkg/pifont
% \newcommand{\cmark}{{\color{ForestGreen}\ding{51}}}%
% \newcommand{\xmark}{{\color{Maroon}\ding{55}}}%

%%%%%
% Comments and TODOs
%%%%%

\usepackage{soul}
\usepackage{hyperref}

\ifDEBUG
    % Floating comment boxes
    \definecolor{olive}{rgb}{0.5, 0.5, 0.0}
    \definecolor{orange}{rgb}{1.0, 0.65, 0.0}
    \newcommand{\JD}[1]{\textcolor{blue}{[Jamie: #1]}}
    \newcommand{\TS}[1]{\textcolor{olive}{[Taylor: #1]}}
    \newcommand{\ST}[1]{\textcolor{red}{[Santiago: #1]}}
    \newcommand{\EB}[1]{\textcolor{orange}{[Ethan: #1]}}
    \newcommand{\KC}[1]{\textcolor{green}{[Kelechi: #1]}}
    \newcommand{\CO}[1]{\textcolor{magenta}{[Chinenye: #1]}}


    % Inline comments and WIP phrasing
    \newcommand{\TODO}[1]{\hl{#1}}
    \newcommand{\CHOP}[1]{}
    \newcommand{\placeholder}[1]{\hl{#1}}


    \hypersetup{
        pdfpagemode=UseNone % don't show thumbnails or bookmarks on opening
    }

\else
    \newcommand{\JD}[1]{}
    \newcommand{\TS}[1]{}
    \newcommand{\ST}[1]{}
    \newcommand{\KC}[1]{}
    \newcommand{\EB}[1]{}
    \newcommand{\TODO}[1]{}
    \newcommand{\CHOP}[1]{}
    \newcommand{\placeholder}[1]{#1}
\fi

%%%%%
% Paragraph-level labels
%%%%%
% Make a \paragraph that is inline and emphasized
\newcommand{\myparagraph}[1]{\vspace{0.08cm}\noindent\hspace{0.1cm} \textit{#1:}}
%\renewcommand{\myparagraph}[1]{\paragraph{#1}}
%\renewcommand{\myparagraph}[1]{\noindent \hspace{0.1cm}\textit{#1}}

\newcommand{\myminisection}[1]{\vspace{0.1cm} \noindent \textbf{#1.}}

\newcommand{\mysubsubsection}[1]{\subsubsection{\textbf{#1}}}

%%%%%%%%
% Peer review
%%%%%%%%


%%%%%%%%
% Balance references
%%%%%%%%

\usepackage{balance} % Put \balance on the last page

%%%%
% Intra-paper references
%%%%
\usepackage{cleveref}

\crefformat{section}{\S#2#1#3}
\crefrangeformat{section}{\S#3#1#4--\S#5#2#6}
\crefname{figure}{Figure}{Figures}
\crefname{appendix}{Appendix}{Appendices}
\crefname{table}{Table}{Tables}
\crefname{algorithm}{Algorithm}{Algorithms}
\crefname{listing}{Listing}{Listings}
\crefname{theorem}{Theorem}{Theorems}
\crefname{thm}{Theorem}{Theorems}
\crefname{lemma}{Lemma}{Lemmata}
\crefname{equation}{Eqt.}{Eqts.}
\crefformat{Grammar}{Grammar #1}


%%%%
% Latin
%%%%

\newcommand{\ie}{\textit{i.e.,} }
\newcommand{\eg}{\textit{e.g.,} }
\newcommand{\etal}{\textit{et al.}\xspace}
\newcommand{\etals}{\textit{et al.'s}\xspace}
\newcommand{\nb}{\textit{N.B.,} }

% Please add the following required packages to your document preamble:

% Nice formatting for a list of RQs
\newenvironment{RQList}{
   \setlength{\topsep}{0pt}
   \setlength{\partopsep}{0pt}
   \setlength{\parskip}{0pt}
   \begin{description}[style=unboxed]
   \setlength{\leftmargin}{1in}
   \setlength{\parsep}{0pt}
   \setlength{\parskip}{0pt}
   \setlength{\itemsep}{0pt}
   }
   {\end{description}}
   
\setlength{\abovecaptionskip}{2pt plus 1pt minus 1pt}
\setlength{\belowcaptionskip}{2pt plus 1pt minus 1pt}

\newcommand{\myquote}[1]{\begin{quote}\emph{``#1''}\end{quote}}
\newcommand{\myinlinequote}[1]{\emph{``#1''}}

\usepackage{quoting} % For customized quotation environments

% Customize the quoting environment
\quotingsetup{
  leftmargin=0.25cm, % Adjust the left margin
  rightmargin=0.25cm, % Adjust the right margin
}

% Define a command to indent the first line
\newcommand{\firstlineindent}{\hspace*{0.3cm}}

\renewcommand{\myquote}[2]{
\begin{quoting}
{
\noindent
\small \firstlineindent \emph{``\textit{#1}'' \textnormal{--- #2}}
}
%\begin{flushright}
%\textnormal{#2}
%\end{flushright}
\end{quoting}}

\newcommand{\SmallLongString}[1]{ {\small \emph{#1}}\xspace}

% ParticipantID, Quote
\newcommand{\ParticipantQuote}[2]{
    \begin{quote}
    \emph{\textbf{#1}: ``#2''} 
    \end{quote}}
    
\DeclareUnicodeCharacter{2060}{\nolinebreak}
